
Parameter-efficient fine-tuning methods like LoRA have enabled task-specific adaptation of large language models with minimal compute. However, their success depends on an implicit assumption: the base model architecture already supports the target task. When this assumption fails, fine-tuning cannot overcome architectural constraints. A state-space model may achieve 81\% accuracy on sentiment classification but perform 12 percentage points below random baseline on paraphrase detection---a failure no parameter tuning can fix.

This work examines when PEFT methods fail due to architectural incompatibility rather than insufficient adaptation. While LoRA and related approaches have been extensively validated on transformer architectures, their applicability to sub-quadratic models like state-space models remains unexplored. Sub-quadratic architectures achieve efficiency through linear-time processing, replacing quadratic attention with recurrent dynamics. However, these architectural changes introduce constraints: causal state-space models process sequences autoregressively without bidirectional context, limiting which tasks they can perform.

The core challenge is architectural compatibility: some tasks require capabilities inherent to the model architecture, independent of learned parameters. Paraphrase detection requires symmetric comparison of two sequences---a capability present in transformer encoders through bidirectional attention but absent in causal decoders. If the base architecture cannot perform a task at zero-shot, parameter-efficient fine-tuning cannot add that capability. PEFT adapts existing representations but cannot alter the information flow dictated by architecture.

Our key insight is that zero-shot evaluation serves as an architectural compatibility gate: if a pretrained model achieves below-random accuracy before fine-tuning, the architectural mismatch signals that PEFT will fail. This validation step prevents wasted compute on structurally impossible experiments.

We evaluate Mamba-130M, a representative causal state-space model, on three GLUE tasks: natural language inference (MNLI), paraphrase detection (QQP), and sentiment classification (SST-2). These tasks probe different reasoning requirements. Zero-shot evaluation reveals a binary pattern: Mamba achieves 81\% on SST-2 (31 percentage points above random, p < 0.001), confirming genuine language understanding. Yet it achieves only 38\% on QQP (12 percentage points below random, p = 0.003), performing worse than guessing. This below-random performance exposes architectural incompatibility that fine-tuning cannot overcome.

We make the following contributions:

\begin{enumerate}
\item \textbf{Systematic evaluation of architectural compatibility for PEFT transfer to state-space models}: We demonstrate that parameter-efficient fine-tuning methods developed for transformers do not automatically apply to causal state-space models when task requirements misalign with architectural capabilities.
\end{enumerate}

\begin{enumerate}
\item \textbf{Empirical demonstration of architectural failure modes}: Mamba-130M fails bidirectional reasoning tasks at zero-shot (QQP: 38\% versus 50\% baseline), confirming that causal architectures cannot perform symmetric comparisons required for paraphrase detection.
\end{enumerate}

\begin{enumerate}
\item \textbf{Validation of zero-shot evaluation as architectural compatibility gate}: Below-random zero-shot accuracy reliably signals PEFT incompatibility, preventing wasted compute on architecturally impossible tasks.
\end{enumerate}

\begin{enumerate}
\item \textbf{Task-architecture compatibility framework}: Architectural capabilities are prerequisites for PEFT, not outcomes. Fine-tuning specializes existing capabilities but cannot add fundamentally new ones.
\end{enumerate}

\textbf{Scope and Limitations:} We evaluate zero-shot architectural compatibility without empirically testing PEFT fine-tuning after observing failure. Our analysis predicts that fine-tuning would amplify architectural bias rather than overcome it, but empirical validation of this prediction remains future work. Results are based on a single state-space model instance (Mamba-130M); broader claims about causal state-space models require testing additional variants. The work focuses on establishing zero-shot evaluation as a gating mechanism rather than comprehensively documenting PEFT failures across all possible configurations.
